\chapter{CHASE: FOLLOWING A DOG OR PERSON}
\label{ch:chase}

The chase behaviour points the robot at a detected dog (or, for testing, a
person) and drives toward it until it is a set distance away. This chapter
describes the whole chain: the on-camera pipeline, the neural network that
recognises the target, how detections become a track, the control law and its
stability, how it shares the motors with teleop, and the plan for bringing it
up on the robot. \S\ref{sec:chase} gives a short summary of the ROS node.

\influx{Chase has only been run in dry-run mode, where it computes commands
but never sends them to the motors. It is waiting on the camera's separate
power supply (\S\ref{sec:powerbudget}).}

\section{PIPELINE}
\label{sec:chasepipe}

\begin{figure}[htbp]
\centering
\begin{tikzpicture}[
  box/.style={draw,rounded corners=2pt,align=center,font=\small,
              minimum height=10mm,inner sep=3pt,text width=24mm},
  arr/.style={-{Stealth[length=2mm]},thick},
  lab/.style={font=\scriptsize,align=center}]
\node[box] (rgb) {IMX214 colour\\1080p, 10 fps};
\node[box,below=6mm of rgb] (mono) {2$\times$ OV7251\\400P, 10 fps};
\node[box,right=8mm of rgb] (prev) {Preview\\$300\times300$};
\node[box,right=8mm of mono] (stereo) {StereoDepth\\aligned to RGB};
\node[box,right=8mm of prev,yshift=-8mm,text width=30mm] (nn)
  {MobileNet-SSD\\spatial detection\\(Myriad X)};
\node[box,right=8mm of nn,text width=26mm] (node) {\texttt{chase\_node}\\select, debounce,\\smooth, control};
\node[box,below=12mm of node,text width=26mm] (base) {\texttt{bbai\_base}\\motors};
\node[box,above=8mm of node,text width=26mm] (joy) {PS4 \topic{/joy}\\R2 arms, L1 overrides};
\draw[arr] (rgb) -- (prev);
\draw[arr] (mono) -- (stereo);
\draw[arr] (prev) -- (nn);
\draw[arr] (stereo) -- (nn);
\draw[arr] (nn) -- node[above,lab]{USB}(node);
\draw[arr] (node) -- node[right,lab]{\topic{/cmd\_vel}\\(armed, not dry run)}(base);
\draw[arr] (joy) -- (node);
\begin{scope}[on background layer]
\node[draw,dashed,rounded corners,fit=(rgb)(mono)(prev)(stereo)(nn),inner sep=4pt,
      label={[font=\scriptsize]above:OAK-D-Lite}] {};
\end{scope}
\end{tikzpicture}
\caption{Chase pipeline. Everything inside the dashed box runs on the
camera; the BeagleBone sees only detections.}
\label{fig:chasepipe}
\end{figure}

Figure~\ref{fig:chasepipe} shows the data flow. The DepthAI pipeline built by
the node has five parts: the colour camera at 1080p producing a
$300\times300$ preview, the two mono cameras at 400P, a \texttt{StereoDepth}
node in high-density mode with depth aligned to the colour camera
(\S\ref{sec:stereo}), and a \texttt{MobileNetSpatialDetectionNetwork} that
runs the detector on the preview and attaches a 3-D position to every
detection (\S\ref{sec:spatial}). The camera runs at the node's
\texttt{fps} parameter, 10 by default, which bounds the detection rate. The
model is the OpenVINO 2021.4 MobileNet-SSD compiled for 6 SHAVE cores
(\file{\textasciitilde/models/mobilenet-ssd\_openvino\_2021.4\_6shave.blob}).

A background thread reads detection messages from the camera and keeps the
latest list of $(\text{label},\text{confidence},X,Z,\bar u)$ tuples, with
$\bar u\in[0,1]$ the box centre across the image. If the camera throws (USB
disconnect, power dip), the thread records the error, reports
\texttt{CAMERA\_ERROR}, and retries every \SI{2}{s}. (On the board,
\texttt{numpy} and \texttt{rospy} must be imported before \texttt{depthai},
or the depthai import fails.)

\section{OBJECT RECOGNITION: MOBILENET-SSD}
\label{sec:mobilenet}

\paragraph{Backbone.} MobileNet \cite{howard2017mobilenets} replaces each
standard convolution with a \emph{depthwise separable} one: a $D_K\times D_K$
convolution applied to each of the $M$ input channels separately, followed
by a $1\times1$ convolution that mixes them into $N$ output channels. On a
$D_F\times D_F$ feature map the cost falls from
$D_K^2MND_F^2$ to $D_K^2MD_F^2+MND_F^2$ multiply-adds, a ratio of
\begin{equation}
\frac{D_K^2MD_F^2+MND_F^2}{D_K^2MND_F^2}=\frac1N+\frac1{D_K^2}\approx\frac19
\quad(D_K=3,\ N\gg9).
\end{equation}
That eight- to nine-fold saving is what lets the detector run at camera
frame rate on the Myriad~X.

\paragraph{Detector head.} SSD \cite{liu2016ssd} predicts objects directly
from several feature maps of decreasing resolution (for MobileNet-SSD at
$300\times300$: $19\times19$, $10\times10$, $5\times5$, $3\times3$,
$2\times2$ and $1\times1$). Each cell of each map carries a few fixed
\emph{default boxes} of different sizes and aspect ratios, 1917 in all. For
every default box $a=(a_x,a_y,a_w,a_h)$ the network outputs 21 class scores
and four offsets $(t_x,t_y,t_w,t_h)$, decoded to a box by
\begin{equation}
b_x=a_x+\sigma_c t_x a_w,\quad
b_y=a_y+\sigma_c t_y a_h,\quad
b_w=a_w e^{\sigma_s t_w},\quad
b_h=a_h e^{\sigma_s t_h},
\end{equation}
with the standard variances $\sigma_c=0.1$, $\sigma_s=0.2$. The class scores
pass through a softmax. Boxes whose best class score is below the confidence
threshold (0.5 here) are dropped, and \emph{non-maximum suppression} removes
duplicates: boxes are taken in descending score, and any box overlapping a
kept box of the same class by intersection-over-union
\begin{equation}
\mathrm{IoU}(A,B)=\frac{\abs{A\cap B}}{\abs{A\cup B}}
\end{equation}
above a threshold (typically 0.45) is discarded. All of this runs on the
camera.

\paragraph{Classes.} The network is trained on PASCAL VOC
\cite{everingham2010pascal}, whose 20 classes plus background include
\emph{dog} (index 12) and \emph{person} (index 15), as well as cat, which
the network may confuse with a small dog. The node's \texttt{targets}
parameter chooses which labels to follow: \texttt{[dog]} by default, and
\texttt{[dog, person]} for testing by walking in front of the camera. In
bench tests the network detected the dog at 93--99\% confidence, about 10
times per second.

\section{FROM DETECTIONS TO A TRACK}
\label{sec:chasemeas}

Each time a new detection frame arrives (about \SI{10}{Hz}) the node:
\begin{enumerate}
\item keeps detections whose label is in \texttt{targets} and takes the most
  confident one;
\item computes the measured bearing and range in the camera frame, using
  depth if it is valid:
\begin{equation}\label{eq:chasemeas}
(\beta_m,\rho_m)=\begin{cases}
\bigl(\atantwo(X,Z),\ \sqrt{X^2+Z^2}\bigr) & Z>\SI{0.05}{m},\\[2pt]
\bigl((\bar u-\tfrac12)\,\theta_h,\ \text{unknown}\bigr) & \text{otherwise},
\end{cases}
\end{equation}
  with $\beta$ positive to the right and $\theta_h$ the parameter
  \texttt{rgb\_hfov\_deg} (see \S\ref{sec:previewcrop} on its value);
\item smooths both with an exponential moving average (EMA) of weight $a=0.5$
  on the new sample,
\begin{equation}\label{eq:ema}
\hat\beta_k = a\beta_{m,k}+(1-a)\hat\beta_{k-1},
\end{equation}
  and the same for $\rho$ (a missing range keeps the previous one);
\item counts consecutive frames with a target, $h$. A frame with no target
  resets $h$ to 0, and the next detection restarts the EMA from the raw
  measurement.
\end{enumerate}
The node publishes the smoothed target as a point on \topic{/chase/target}
in the camera frame.

\paragraph{The EMA as a filter.} With samples every $T=\SI{0.1}{s}$,
\eqref{eq:ema} is the discretization of a first-order low-pass filter with
time constant
\begin{equation}
\tau_e=\frac{-T}{\ln(1-a)}=\frac{0.1}{\ln2}\approx\SI{0.14}{s},
\end{equation}
which cuts the variance of white bearing noise to a third
($\operatorname{var}\hat\beta = \frac{a}{2-a}\operatorname{var}\beta_m=\frac13\operatorname{var}\beta_m$)
at the cost of \SI{0.14}{s} of lag.

\paragraph{Debounce and loss.} A command is produced only when $h\ge3$ (three
consecutive frames, \SI{0.3}{s}) and the last detection is less than
\SI{0.5}{s} old. The three-frame rule rejects one-off false positives, but
because a single missed frame resets $h$, it also means one dropped detection
stops the robot for at least \SI{0.3}{s}. Table~\ref{tab:chasestate} lists the
states reported on \topic{/chase/status} twice a second.

\begin{table}[htbp]
\centering
\caption{Chase states}
\label{tab:chasestate}
\small
\begin{tabularx}{\linewidth}{lX}
\toprule
State & Meaning \\
\midrule
\texttt{SEARCH} & No usable target: none seen, fewer than 3 consecutive hits, or the last one older than \SI{0.5}{s}. Commands zero. \\
\texttt{TRACK} & Target held; the status line shows label, confidence, range, bearing and the command. \\
\texttt{CAMERA\_ERROR} & The camera thread failed; it retries every \SI{2}{s}. \\
\texttt{ARMED} / \texttt{idle} & Suffix: whether R2 is held (and L1 is not). \\
\texttt{(dry run)} & Suffix: commands go only to \topic{/chase/cmd\_vel\_preview}. \\
\bottomrule
\end{tabularx}
\end{table}

\section{CONTROL LAW}
\label{sec:chaselaw}

At \SI{15}{Hz} the node turns the track into a velocity command:
\begin{align}
\omega &= \begin{cases}
  \sat_{\omega_{\max}}\!\bigl(-k_\beta\hat\beta\bigr) & \abs{\hat\beta}>\beta_{\mathrm{db}},\\
  0 & \text{otherwise},\end{cases}\label{eq:chasew}\\
v &= \sat_{[0,v_{\max}]}\!\bigl(k_\rho(\hat\rho-\rho_f)\bigr)\cdot
     \sat_{[0,1]}\!\Bigl(1-\frac{\abs{\hat\beta}}{\beta_{\mathrm{full}}}\Bigr)
     \quad\text{if }\hat\rho>\rho_{\mathrm{stop}},\ \text{else }0,\label{eq:chasev}
\end{align}
with the parameters in Table~\ref{tab:chaseparams}. The minus sign in
\eqref{eq:chasew} is because a target to the right ($\beta>0$) needs a
clockwise turn, which is negative $\omega$ in REP-103. The second factor in
\eqref{eq:chasev} scales forward speed down while the robot is still turning
toward the target, reaching zero at \SI{25}{\degree} off-axis, so the robot
turns first and then drives. $v$ is never negative: the robot does not back
away from a target that comes too close, it just stops.

\begin{table}[htbp]
\centering
\caption{Chase parameters (\texttt{chase.launch})}
\label{tab:chaseparams}
\small
\begin{tabular}{llll}
\toprule
Symbol & Parameter & Value & Role \\
\midrule
$k_\beta$ & \texttt{k\_angular} & \SI{2.5}{s^{-1}} & turn rate per radian of bearing \\
$\omega_{\max}$ & \texttt{max\_angular} & \SI{2.0}{rad/s} & turn-rate limit \\
$\beta_{\mathrm{db}}$ & \texttt{bearing\_deadband\_deg} & \SI{5}{\degree} & no turning inside this \\
$\beta_{\mathrm{full}}$ & \texttt{full\_speed\_bearing\_deg} & \SI{25}{\degree} & forward speed zero beyond this \\
$k_\rho$ & \texttt{k\_linear} & \SI{0.4}{s^{-1}} & forward speed per metre of range error \\
$v_{\max}$ & \texttt{max\_linear} & \SI{0.2}{m/s} & forward speed limit (deliberately slow) \\
$\rho_f$ & \texttt{follow\_distance} & \SI{0.9}{m} & range it settles at \\
$\rho_{\mathrm{stop}}$ & \texttt{stop\_distance} & \SI{0.5}{m} & no forward motion inside this \\
$a$ & \texttt{smoothing} & 0.5 & EMA weight \\
 & \texttt{min\_hits}, \texttt{lost\_timeout} & 3, \SI{0.5}{s} & debounce and loss \\
 & \texttt{rate}, \texttt{fps} & \SI{15}{Hz}, 10 & control and camera rates \\
\bottomrule
\end{tabular}
\end{table}

\subsection{Bearing loop}

For a stationary target, turning the robot left by $\dd\psi$ moves the
target right by the same angle, so $\dot\beta=\omega_{\mathrm{actual}}$.
Ignoring saturation and the dead band, and if the robot turns at the
commanded rate, the loop is
\begin{equation}
\dot\beta = -k_\beta\,\beta,
\end{equation}
which converges exponentially with time constant $1/k_\beta=\SI{0.4}{s}$.

\paragraph{Delay and stability.} The real loop has delay: the camera frame
period and network latency, the \SI{15}{Hz} control and \SI{30}{Hz} base
loops, and the motor's response, which together are roughly
$\tau_d\approx\SI{0.2}{s}$, plus the EMA's lag $\tau_e=\SI{0.14}{s}$. The
open-loop transfer function is then
\begin{equation}
L(s)=\frac{k_\beta\,e^{-s\tau_d}}{s(\tau_e s+1)} .
\end{equation}
Its gain crosses one at $\omega_c\approx\SI{2.4}{rad/s}$, where the phase
margin is
\begin{equation}
\phi_m = \SI{90}{\degree}-\frac{\SI{180}{\degree}}{\pi}\,\omega_c\tau_d-\arctan(\omega_c\tau_e)
\approx 90-28-19=\SI{44}{\degree}.
\end{equation}
That is adequate: the step response will overshoot slightly but settle. The
loop would go unstable at a total delay of about \SI{0.52}{s}, and the
margin shrinks quickly with gain (Table~\ref{tab:chasepm}), which is why
\emph{raising $k_\beta$ is not the way to make the robot turn harder.}

\begin{table}[htbp]
\centering
\caption{Bearing-loop phase margin, $\tau_e=\SI{0.14}{s}$}
\label{tab:chasepm}
\small
\begin{tabular}{lccc}
\toprule
 & $\tau_d=\SI{0.15}{s}$ & $\tau_d=\SI{0.2}{s}$ & $\tau_d=\SI{0.3}{s}$ \\
\midrule
$k_\beta=2.5$ & \SI{51}{\degree} & \SI{44}{\degree} & \SI{30}{\degree} \\
$k_\beta=4$ & \SI{32}{\degree} & \SI{22}{\degree} & \SI{2}{\degree} \\
$k_\beta=6$ & \SI{13}{\degree} & unstable & unstable \\
\bottomrule
\end{tabular}
\end{table}

\paragraph{Will the robot actually turn?} The bearing loop assumes the
robot turns at the commanded rate, and on this robot it does not
(\S\ref{sec:skidturn}). The detector only sees the central part of the
colour frame, about $\pm\SI{21}{\degree}$ if the preview is cropped
(\S\ref{sec:previewcrop}), so the largest command chase will ever send is
$2.5\times0.37\approx\SI{0.92}{rad/s}$, and even the \SI{2}{rad/s} limit maps
through \eqref{eq:sidespeed}--\eqref{eq:duty} to a duty of only 0.16--0.21.
Teleop's old \SI{1.5}{rad/s} full stick (duty 0.19) could not start a spin
on the floor. Chase will therefore very likely stall on turns: the motors
form a dead zone in the loop, the bearing error sits there, and nothing
happens until $v$ (which is held near zero by the bearing factor) or a
nudge breaks it free. The fix belongs below the chase node, in the yaw-rate
loop of \S\ref{sec:yawloop}, so that a commanded rate is delivered on any
floor; until then a larger turning duty floor in \texttt{bbai\_base} for
spins in place is the quick workaround.

\subsection{Range loop}

For a stationary target straight ahead, $\dot\rho=-v$, and
\eqref{eq:chasev} gives $\dot\rho=-k_\rho(\rho-\rho_f)$ away from the limits:
the range settles at $\rho_f$ with time constant $1/k_\rho=\SI{2.5}{s}$.
Beyond $\rho_f+v_{\max}/k_\rho=\SI{1.4}{m}$ the speed saturates and the robot
closes at a steady \SI{0.2}{m/s}. A target receding at speed $u<v_{\max}$ is
followed with a steady lag $\rho-\rho_f=u/k_\rho$ (\SI{0.25}{m} at
\SI{0.1}{m/s}); a target faster than \SI{0.2}{m/s}, which includes a walking
dog, is lost. At these settings the behaviour is a slow \emph{follow}, not a
chase, which is the intent for first tests. The range loop needs no delay
analysis: its bandwidth is far below the delay.

\section{ARBITRATION AND SAFETY}
\label{sec:chasesafety}

\begin{itemize}
\item \textbf{Dry run by default.} Unless launched with
  \texttt{dry\_run:=false}, chase never publishes \topic{/cmd\_vel}; it only
  publishes what it would send on \topic{/chase/cmd\_vel\_preview}.
\item \textbf{Dead-man arming.} Commands go to \topic{/cmd\_vel} only while
  R2 (button 7) is held and L1 (button 4, teleop's dead-man) is not. Releasing
  R2 sends one zero command. Teleop and chase therefore never publish at the
  same time, and the base node, which applies whichever command arrived last,
  needs no arbitration logic of its own.
\item \textbf{Timeouts.} Chase stops if \topic{/joy} is silent for
  \SI{0.5}{s} (controller out of range) or the target is lost for
  \SI{0.5}{s}; the base node stops the motors if \topic{/cmd\_vel} is silent
  for \SI{0.5}{s}. Every failure ends in stopped motors.
\item \textbf{Speed.} \SI{0.2}{m/s} forward, no reverse, no forward motion
  closer than \SI{0.5}{m}.
\item \textbf{No obstacle avoidance.} Chase drives straight at the target and
  does not look at the lidar. Anything between the robot and the dog will be
  hit at up to \SI{0.2}{m/s}.
\end{itemize}

\section{BRING-UP PLAN}
\label{sec:chaseplan}

Once the camera runs from the battery bank through the USB splitter:
\begin{enumerate}
\item \textbf{Power soak.} Run \texttt{robot.launch}, lidar and chase (dry
  run) together for several minutes while driving with teleop, and confirm no
  reset and no \texttt{CAMERA\_ERROR}.
\item \textbf{Dry run with a person.} \texttt{dry\_run\_test.sh} launches chase
  with \texttt{targets:=[dog, person]}, then prints the status, the preview
  commands and the publishers of \topic{/cmd\_vel}, which must not include
  chase. Walk across the field of view and check that the sign of
  \texttt{angular.z} is positive when you are on the robot's left.
\item \textbf{Check the field of view.} Note the bearing at which detections
  stop at the frame edges; about \SI{21}{\degree} confirms the cropped preview
  (\S\ref{sec:previewcrop}).
\item \textbf{Wheels lifted, armed.} Launch with \texttt{dry\_run:=false},
  hold R2, and check that the wheels turn the right way as a person moves.
\item \textbf{On the floor with a person.} Check turning on the actual floor
  (the stall described above), the follow distance and the stop.
\item \textbf{With the dog.} Only after the above, with \texttt{targets:=[dog]}.
\end{enumerate}

\section{PROPOSED IMPROVEMENTS}
\label{sec:chasefuture}

\emph{None of these is implemented.}

\paragraph{Hysteresis on loss.} Decrement the hit counter on a missed frame
instead of resetting it, and keep commanding the last smoothed bearing until
the \SI{0.5}{s} loss timeout, so a single dropped detection does not stop the
robot.

\paragraph{Target filter in the odometry frame.}
\label{sec:targetkf}
Today the bearing mixes the target's motion with the robot's own rotation,
and the track is lost the moment the target leaves the image. Tracking the
target in \frm{odom} separates the two. With target state
$\bm s=(p_x,p_y,\dot p_x,\dot p_y)$ and a constant-velocity model,
\begin{equation}
\bm s_{k+1}=\begin{bmatrix}I&T I\\0&I\end{bmatrix}\bm s_k+\bm w_k,\qquad
\bm w_k\sim\mathcal N\!\left(0,\ q\begin{bmatrix}\tfrac{T^3}{3}I&\tfrac{T^2}{2}I\\
\tfrac{T^2}{2}I&TI\end{bmatrix}\right),
\end{equation}
each detection is mapped into \frm{odom} through the camera transform
\eqref{eq:cam2base} and the robot pose $(x,y,\psi)$,
\begin{equation}
\bm z_k=\begin{bmatrix}x\\y\end{bmatrix}+R(\psi)\left(
\begin{bmatrix}Z\\-X\end{bmatrix}+\bm t_{\mathrm{cam}}\right),
\qquad
R_z = J\,\operatorname{diag}\!\bigl(\sigma_\beta^2\rho^2,\ \delta Z^2\bigr)J\T,
\end{equation}
where $J$ rotates the bearing/range errors into \frm{odom} and $\delta Z$
comes from \eqref{eq:depthres}. A standard Kalman filter
(Chapter~\ref{ch:fusion}) then gives a target position and velocity that
survive short dropouts, can be extrapolated over the detection delay, and
can steer the robot back toward where the target left the image.

\paragraph{Search.} When the target has been lost for longer than a few
seconds, turn slowly in place toward its last known side. This needs reliable
turning first (\S\ref{sec:yawloop}).

\paragraph{Obstacle gating.} Limit forward speed from the lidar's closest
return $d$ in a sector ahead of the robot,
$v\le\sqrt{2a_b\,\max(0,d-d_s)}$, with a braking deceleration $a_b$ and a
safety distance $d_s$, so the robot can always stop before an obstacle. A
full alternative is to send the target position as a \texttt{move\_base}
goal once navigation is set up.

\paragraph{Faster detector input.} Running the camera at 15--20 fps
reduces the delay and allows a higher $k_\beta$ for the same margin, at
the cost of more camera power.
